\chapter{PART V · 技能重构}


\newpage

技能是组织的血液。

技能决定了组织能做什么、不能做什么。当技术改变时，技能的需求也在改变——有些技能变得更有价值，有些技能变得不那么重要，有些技能需要被重新定义，有些技能需要被全新创造。

AI时代，技能的重构正在以前所未有的速度发生。机器正在学会做那些过去只有人类才能做的事。这意味着，那些过去定义"好员工"的技能标准，正在被重新改写。

这一部分要探讨的，是AI时代人才发展需要关注的能力维度。核心命题是：\textbf{可编码的技能在贬值，不可编码的技能在增值。} 那些难以被标准化的能力——判断力、创造力、关系力——变得比以往任何时候都更稀缺、更值钱。

五个章节，对应技能重构的五个关键维度：
\begin{itemize}
\item 技能经济学的重构——可编码与不可编码的此消彼长
\item T型能力的进化——从深度+广度到整合深度
\item 判断力：人类最后的护城河——为什么判断力在AI时代变得更重要
\item 人机协作技能——与AI共事的新基本功
\item 关系技能：管理的本质——为什么关系力在AI时代变得更稀缺

\end{itemize}
这个重构的核心，是重新定义"人才"的含义——在AI可以完成大多数重复性、程序性任务的世界里，人类的独特价值是什么？
